\documentclass[letter, 10pts]{article}
\usepackage[monocolor]{../math232/ahsansabit}
\usepackage[]{float}
\usepackage{tikz}
\usepackage{tikz-3dplot}
\usepackage[outline]{contour} % glow around text
\usepackage{xcolor}
\usepackage{pdfpages}
\usepackage{physics}
\usepackage{multicol}
\title{Quantum Mechanics : : Homework 08}
\author{Ahmed Saad Sabit, Rice University}
\date{\today}
\newcommand{\hb}{\hbar}
\newcommand{\U}{\uparrow}
\newcommand{\D}{\downarrow}
\usepackage[]{braket}
\begin{document}
\maketitle

Extension granted. 

\section*{Problem 01}
\hrule
Used in Integral Calculator: 
\[
	\psi_0(x) = \left(\frac{ m \omega}{ \pi \hb } \right) ^{ 1 / 4} e^{ - m \omega x^2 /  2 \hb }
\]
For $n = 1$ eigenstate wavefunction keeping a common factor between two wavefunctions
\[
	\psi_1(x) = x \left(\frac{ m \omega}{ \pi \hb } \right) ^{ 1 / 4} \sqrt{\frac{2 m \omega }{\hb } } e^{ - m \omega x^2 /  2 \hb }
\] 
\subsection*{(a)}
Without $C \delta(x_1 - x_2)$ repulsion, the two particles could happily have the same wave functions $\psi(x) \sigma_\text{spin}$. Without the perturbation 
\[
\hat{H} \ket{\psi_1 , \psi_2} = \hat{H}_1 \ket{\psi_1 \psi_2} + \hat{H}_2 \ket{\psi_1\psi_2}  \implies \varepsilon \ket{\psi_1 \psi_2}
\]
\[
E_0 ^{[0] } = 2 \left(\frac{\hb \omega }{2 }\right) = \hb \omega 
\] 
With the perturbation
\[
\hat{H} \ket{\psi_1 , \psi_2} = \hat{H}_1 \ket{\psi_1 \psi_2} + \hat{H}_2 \ket{\psi_1\psi_2} + C \delta(\hat x_1 - \hat x_2) \ket{\psi_1 \psi_2} \implies E \ket{\psi_1 \psi_2}
\]
The energy level perturbation of first order


\[
\Delta = \braket{\psi_1 \psi_2 | C \delta( \hat x_1 -  \hat{ x }_2) | \psi_1 \psi_2} = C
\int \Psi(x_1, x_1) ^{*} \Psi(x_1, x_1) \, \mathrm{d} x_1
\] 

Unperturbed lowest energy state (using integral calculator)
\[
C \int_{-\infty}^{\infty} \mathrm{d} x \, | \Psi(x) | ^{4}
= 
C \sqrt{\frac{m \omega}{2 \pi \hb}} 
\] 

\[
\boxed{
	E_0 ^{[1]} = \hb \omega + C \sqrt{\frac{m \omega}{2 \pi \hb}} 
}
\] 



\subsection*{(b)} 
Energy without perturbation
\[
E = \frac{1}{2} \hb \omega + \frac{3}{2} \hb \omega = 2 \hb \omega 
\] 
Symmetric orbital is given by 
\[
\psi(x_1 , x_2) _s = \frac{1}{\sqrt{2} } 
\left(
\psi_0(x_1) \psi_1(x_2) + \psi_0 (x_2) \psi_1 (x_1)
\right)
\] 
Anti-Symmetric orbital is given by 
\[
\psi(x_1 , x_2) _A = \frac{1}{\sqrt{2} } 
\left(
\psi_0(x_1) \psi_1(x_2) - \psi_0 (x_2) \psi_1 (x_1)
\right)
\] 

It's easy compute the energy for anti-symmetric wavefunction because $\psi(x_1, x_2)_A = 0$ for $x_1 = x_2$ giving $E_A ^{(1) } = 0$. The orbital wavefunction is anti-symmetric for spin triplet state. 

For symmetric orbital wavefunction 
\[
E_S ^{(1)} = 
\int_{-\infty}^{\infty}  
\int_{-\infty}^{\infty}  
\mathrm{d} x_1 \, \mathrm{d} x_2 \, 
\psi_S^{*} (x_1 ,x_2) C \delta(x_1  - x_2) \psi_S(x_1 , x_2) = \int_{-\infty}^{\infty}  \mathrm{d}  x \, 2 C \psi_0 ^2 (x) \psi_1 ^2 (x) = C \sqrt{\frac{m \omega }{2 \pi \hb}} 
\] 


The 4 degeneracy of the energy has been converted to 3 degeneracy in Anti-symmetric orbital wave function (triplet spin state) and there's only 1 state in single case with a different energy given (both) by
\[
\boxed{
E_S ^{(1)} = C \sqrt{\frac{m \omega}{2 \pi \hb}}  \quad \text{and} \quad E_A ^{(1)} = 0
}
\] 
\section*{Problem 02}
\subsection*{(a)}
Plugging in $K = 1$ 
hence solving the equation for expectation 
value is with $Z_e = 2 $
\[
2 * 2 \frac{1 + 3 + 1}{2^3	 }  (13.6 \, \text{eV}) = 34.0 \text{ eV}
\] 

\subsection*{(b)} 
\[
	\braket{H } = 
	\braket{\psi_{100} \psi_{100}  | \frac{p^2}{2m} + \frac{P^2}{2 M} - 
	\frac{2 e^2}{R} - \frac{2 e^2}{r} + \frac{e^2}{|\vec{r} - \vec{R}| } | \psi_{100} \psi_{100} }
\] 
\[
= E_\text{electron ground} + E_\text{muonic ground} + E_\text{interaction}
\]

Using $11.66$ and $11.72 $ in Townsend we can find 
\[
\braket{nlm | { \hat{p}^2/2 \mu }| nlm} = - E_n ^{(0)} = \frac{\mu c^2 Z^2 \alpha^2}{2 n^2} 
\implies 
E_n ^{(0)}\approx
 - \frac{m_e c^2 Z^2 \alpha ^2}{2 n ^2}
\] 

\[
- \frac{1}{2} \braket{nlm | \frac{ Z e^2 }{\hat{r}} | n l m} =  E_n^{(0)} \implies 
\braket{nlm | \frac{ Z e^2 }{\hat{r}} | n l m} =  -2 E_n ^{(0)} \approx  \frac{m_e c^2 Z^2 \alpha ^2 }{n^2 }
\] 

Invoke $R_\infty$ then 
\[
E_n^{(0)} = - \frac{Z ^2 R_\infty}{n^2} = - \frac{1}{n^2} \left(
\frac{m_e c^2 Z^2 \alpha ^2}{ 2}
\right)
\] 

\[
  - 2 E_n ^{(0) } =  \frac{2 Z  R_\infty}{n^2}
\] 

So for electron $E_n ^{(0)}$ 


Total kinetic energy term for electron and muon 
\[
E_k^{(0)} = R_\infty \left(Z_e ^2 + Z_\mu ^2 \frac{M}{m}\right)
\] 
Total potential energy
\[
E_p^{(0)} = -
4 R_\infty \left(Z_e  + Z_\mu  \frac{M}{m}\right)
\] 

Total Energy 
\[ E = 
 R_\infty \left(Z_e ^2 + Z_\mu ^2 \frac{M}{m}\right)
 - 4 R_\infty \left(Z_e  + Z_\mu  \frac{M}{m}\right) 
 + 
 2 Z_e \frac{(1 + 3 K + K ^2 ) }{(1 + K)^
 3} R_\infty
\]

\subsection*{(c)} 


Because of the following fact
\[
\left|  
\frac{d}{d K} 
\frac{1 + 3 K + K^2}{(1 + K ) ^3 } \right|_{K = 0 }  = 0
\] 
We can make the approximation 

\[
0 = \frac{\partial E}{\partial Z_e} = 2 R_\infty Z_e - 4 R_\infty + 2 R_\infty \implies Z_e = 1 
\]

\[
0 = \frac{\partial E}{\partial Z_\mu} = 2 Z_\mu R_\infty \frac{M}{m} - 4 R_\infty \frac{M}{m} \implies Z_\mu = 2
\] 

\[
K = \frac{1}{2} \frac{m}{207 m} =\frac{1}{414} \implies \frac{1 + 3 K + K^2}{ (1 + K )^3 } \approx
1 \] 

\[
E_\text{ground state} = - 829 R_\infty
\] 


\subsection*{(d)} 
Given the higher mass of muon than electron, we should expect the muon to be much closer to the nucleus than the electron (in a somewhat semi-quasi-classical preview). The effective charge hence shifts and the electron witnesses $Z_e  = 1 $ and $Z_\mu = 2$






\section*{Problem 03}
\hrule 

\subsection*{(a)} 
\[
1s ^2 2s ^2 2 p ^{6} 3 s ^2 3 p ^3
\]
We can count $3$ unpaired valence electron which means total spin is $\frac{3}{2} $ giving us $2 s + 1 = 4$. 

$
L_\text{tot} = | -1 + 0 + 1 | = 0 
$. The three valence electrons are at different states so $J = | L - S | = \frac{3}{2}$. The shell being half full we know $^{4} S _{3 / 2}$


\subsection*{(b)} 
\[
1 s ^2 2 s ^2 2 p ^{6} 3 s ^2 3 p ^{4} 
\]
\[
S_\text{tot} = \text{(2 in pair, 2 without)  } = 1 \implies 2 S + 1 = 3
\]
\[
L_\text{tot} = | - 1 + 0 + 1 \pm 1 | = 1  \text{ which is $P$}
\]
\[
J  = L + S = 2 \implies ^{3} P_2
\] 


\subsection*{(c)} 
\[
1 s ^2 2 s ^2 2 p ^{6} 3 s ^2 3 p ^{6} 4 s ^2 3 d ^{3}
\]
\[
S_\text{tot} = \text{(3 without pairing)  } = \frac{3}{2} \implies 2 S + 1 = 4
\]
\[
L_\text{tot} = | - 2 - 1 + 0 + 1 + 2 - 2 - 1 + 0 | = 3  \text{ which is $F$}
\]
\[
	J  = \text{3 unpaired hence} = | J - S |  = \frac{3}{2} \implies ^{4} F_{3 / 2}
\]


\subsection*{(d)} 
\[
1 s ^2 2 s ^2 2 p ^{6} 3 s ^2 3 p ^{6} 4 s ^2 3 d ^{8}
\]
\[
S_\text{tot} = \text{(2 without pair)  } = 1 \implies 2 S + 1 = 3
\]
\[
L_\text{tot} = | - 2 - 1 + 0 + 1 + 2 - 2 - 1 + 0 | = 3  \text{ which is $F$}
\]
\[
J  = L + S = 4 \implies ^{3} F_4
\]



































\end{document}
